\section{Compatibility of the prime selection and the cohomological interval}
\label{sec:selection-assembly}

This section checks the order of choices in the arithmetic construction.
The numerical selection is supplied unconditionally by
Theorem~\ref{thm:unconditional-qs}, and the common-interval theorem is
used under its stated coarse cohomological hypothesis. The stronger
all-height source zero-free clause is not needed in this partial assembly. In particular, the
common-interval theorem controls the prime $2$ and specified fixed
odd-prime fields; the estimates for all the remaining character components
are identified separately below.

\subsection{The fixed field and the growing cyclotomic extension}

Fix a primitive Dirichlet character $\psi$ of conductor $M\ge1$ and
order $o$. Let $K_\psi$ be the fixed field of $\ker\psi$ in
$\mathbb Q(\mu_M)$. Thus $[K_\psi:\mathbb Q]=o$. Put
\[
 R_h=\mathbb Q(\mu_{2^{h+2}})^+,
 \qquad B_h=K_\psi R_h,
 \qquad E_h=B_h(\sqrt{-q},\sqrt{-d}),
\]
where $q,d$ are distinct primes prime to $2M$, both congruent to $3$
modulo $4$. The superscript $+$ denotes the maximal real subfield.

\begin{lemma}\label{sa-fields}
Write $a=v_2(M)$ and $h_*=\max(a-2,0)$. There is an integer
$1\le c\le o$, depending only on $\psi$, such that, for $h\ge h_*$,
\[
 [B_h:\mathbb Q]=c2^h,\qquad [E_h:\mathbb Q]=4c2^h.
\]
One may take
$c=[K_\psi:K_\psi\cap\bigcup_hR_h]$.
The field $E_h$ is abelian and CM, and is contained in
$\mathbb Q(\mu_{N_h})$, where
\[
 Q_h=\operatorname{lcm}(2^{h+2},M),\qquad N_h=Q_hqd.
\]
In any cyclotomic overfield containing these roots of unity, the
automorphism $\zeta\mapsto\zeta^u$ fixes $E_h$ if and only if
\begin{equation}\label{sa-field-kernel}
 u\equiv\pm1\pmod{2^{h+2}},\quad
 \psi(u)=1,\quad (u/q)=(u/d)=1.
\end{equation}
Here $u$ is prime to the conductor of the overfield.
\end{lemma}

\begin{proof}
Write $M=2^am$ with $m$ odd. The Chinese remainder theorem for
cyclotomic Galois groups shows, for $n\ge a$, that
$\mathbb Q(\mu_M)\cap\mathbb Q(\mu_{2^n})=
\mathbb Q(\mu_{2^a})$; the factors of odd and $2$-power conductor
are independent. Consequently $K_\psi\cap\bigcup_hR_h$ is
already contained in $R_h$ when $h+2\ge a$. The compositum degree
formula and $[R_h:\mathbb Q]=2^h$ give the first assertion.
The integer $c$ is a relative degree and is at most $o$.

For completeness, the quadratic subfield of $\mathbb Q(\mu_q)$
is $\mathbb Q(\sqrt{-q})$ in the present congruence class. Indeed,
for $\zeta=\exp(2\pi i/q)$ and
$\tau=\sum_{u\bmod q}(u/q)\zeta^u$, changing variables gives
$\sigma_v\tau=(v/q)\tau$ and $\overline\tau=(-1/q)\tau$.
Also, writing $u=tv$ in $\tau\overline\tau$, the inner sum over
$v\ne0$ is $q-1$ if $t=1$ and $-1$ otherwise. Hence
$\tau\overline\tau=q$, and $\tau^2=(-1/q)q=-q$.
The same proof applies to $d$. These identities and Galois
correspondence prove the containment and \eqref{sa-field-kernel}.
They also identify any field specified by that fixed-subgroup
condition with $E_h$ up to a $\mathbb Q$-isomorphism.

The field $B_h$ is unramified at $q,d$: at a prime not dividing
its cyclotomic modulus, the polynomial $X^n-1$ has invertible
derivative at every root, so its cyclotomic factors have
discriminant prime to that prime. The three quadratic
subfields of $\mathbb Q(\sqrt{-q},\sqrt{-d})$ have discriminants
$-q,-d,qd$, respectively, so none is contained in $B_h$.
The two fields are therefore linearly disjoint: a nontrivial
intersection with a biquadratic extension would contain a
quadratic subfield. This proves the factor $4$ in the degree.
Finally $E_h$ is abelian and contains the imaginary quadratic
field $\mathbb Q(\sqrt{-q})$. Complex conjugation has order two,
and its fixed field is totally real, so $E_h$ is CM.
\end{proof}

The displayed $Q_h$ is a permissible comparison modulus. If a
construction uses a modulus $Q$ with
$2^{h+2}\mid Q$, $M\mid Q$, and $Q\mid Q_h$, then necessarily
$Q=Q_h$. In particular its odd part is the fixed odd part of $M$.
There is no need to select a new family of base characters modulo
the growing power of $2$.

\begin{lemma}\label{sa-character-families}
Characters of $\operatorname{Gal}(E_h/\mathbb Q)$ trivial on
$\operatorname{Gal}(E_h/K_\psi(\sqrt{-q},\sqrt{-d}))$ are exactly
\[
             \psi^r\chi_q^i\chi_d^j,
 \qquad 0\le r<o,\quad i,j\in\{0,1\}.
\]
They are understood through their primitive Dirichlet inducers,
and $\chi_q(n)=(n/q)$, $\chi_d(n)=(n/d)$.
The family with $(i,j)\ne(0,0)$ has $3o$ distinct members.
Every primitive real character of conductor dividing $Q_hqd$
has the form $\rho\chi_q^i\chi_d^j$, where $\rho$ is primitive
real and $\operatorname{cond}\rho\mid8M$.
\end{lemma}

\begin{proof}
The preceding disjointness gives
$\operatorname{Gal}(K_\psi(\sqrt{-q},\sqrt{-d})/\mathbb Q)
\simeq\operatorname{Gal}(K_\psi/\mathbb Q)\times C_2\times C_2$.
The first factor is cyclic with character group generated by $\psi$.
Taking characters proves the first assertion and distinctness.
If the primitive conductor of $\psi^r$ is $f_r\mid M$, the primitive
conductor after twisting is exactly $f_rq^id^j$, because the three
conductors are coprime. The local factors at omitted primes are
therefore not silently retained in a primitive sum.

For the real-character assertion, a homomorphism from the units
modulo an odd prime power $p^b$ to $\{\pm1\}$ kills the subgroup
of order $p^{b-1}$ and factors modulo $p$. At $2$, every unit
congruent to $1$ modulo $8$ is a square modulo $2^b$. This follows
by induction: if an odd $v$ solves $v^2\equiv u\pmod{2^n}$,
$n\ge3$, replacing $v$ by $v+2^{n-1}$ toggles the next binary
digit of its square. Thus a real character at $2$ factors modulo
$8$. Apply the Chinese remainder theorem to $Q_hqd$. At $q,d$
the real local character is either trivial or the Legendre symbol;
the remaining primitive conductor divides $8M$.
\end{proof}

The characters requiring the square-root estimate in the construction
are among the first family with $(i,j)\ne(0,0)$: the untwisted powers
of $\psi$ are the retained characters. The real-character assertion
also verifies that the zero-free-box requirement involves a fixed
base family. Its primitive conductor is at most
$8M\exp(z+z')$ when $q\le e^z$, $d\le e^{z'}$.
Theorem~\ref{qs-selection} supplies these boxes at every height under
its stated QRH assumption. For the numerical use in this section,
the unconditional Theorem~\ref{thm:unconditional-qs} instead supplies
the finite-height boxes and the required thin bound.

\subsection{Size bounds with the selection scale kept fixed}

Put $\alpha=2+8\epsilon$, where $0<\epsilon<1/52$.
Fix $J\ge1$ and write $U=X^{1/J}$ and $z=U^{\epsilon/5}$.
The prime selection is made at $U$, not at the later variable $x$.
For the selected primes one has
\begin{equation}\label{sa-prime-ranges}
 \frac z2<\log q\le z,\qquad
 \frac{z(\log U)^2}{2}<\log d\le z(\log U)^2.
\end{equation}
For all sufficiently large $X$, these imply
\begin{equation}\label{sa-size-comparison}
 \log(qd)\le X^{9\epsilon/(40J)},\qquad
 q,d>2X^{401}+1,
\end{equation}
and $q,d$ exceed any prescribed fixed constant.
Indeed $\log(qd)\le U^{\epsilon/5}(1+(\log U)^2)$, while
$9/40-1/5=1/40$. For the lower bounds,
$q,d>\exp(X^{\epsilon/(5J)}/2)$ once $\log U\ge1$;
this exceeds every fixed power of $X$.
One explicit verification uses
$\delta=\epsilon/(5J)$,
$n=\lceil402/\delta\rceil$ and $T=2^nn!$:
\[
 \exp(X^\delta/2)\ge X^{\delta n}/T\ge X^{402}/T.
\]
Taking $X>T(B+3)$ and $X\ge e^J$ gives all the bounds above
with an additional prescribed bound $B\ge1$.

If $U\le x\le X$ and $x^\alpha\le c2^h\le2x^\alpha$,
then $Q_h\le8Mx^\alpha$. Consequently, uniformly in such $x,h$,
\begin{equation}\label{sa-conductor}
 \log N_h\le2U^{\epsilon/5}(\log U)^2,
 \qquad
 \log(N_ht)\le2x^{\epsilon/4}\quad(1\le t\le x^{401})
\end{equation}
for sufficiently large $X$. For the first inequality, use
$\log x\le J\log U$ and absorb
$\log(8M)+\alpha J\log U+z$ into $z(\log U)^2$.
For the second use \eqref{sa-size-comparison}, $U\le x$, and
$9\epsilon/40<\epsilon/4$; the remaining logarithm is at most
$\log(8M)+(\alpha+401)\log x$.
Every threshold here depends only on the fixed data.

The congruences are also compatible without additional conditions:
$d\equiv7\pmod{16}$ and $d\equiv-1$ modulo $q$ and each odd
prime at most $z$ are reduced classes at pairwise coprime moduli.
For large $U$, $z\ge4M$ and $q>3$. Each odd prime $\ell\le z$
and the prime $q$ split in $\mathbb Q(\sqrt{-d})$, since
$-d\equiv1\pmod\ell$. The prime $2$ splits as well, since
$-d\equiv1\pmod8$. Requiring $d\equiv-1$ modulo higher prime
powers dividing $M$ is unnecessary for these splitting conditions.

\subsection{A simultaneous interval with the exact rescaling}

\begin{proposition}\label{sa-compatible-interval}
Assume \eqref{ct-coarse-hypothesis}.
Fix $\psi$, a finite set $P$ of odd primes, and fixed abelian fields
$K_p\supset\mu_p$ for $p\in P$. Choose the generator bounds
$R_2,R_p\ge1$ of Proposition~\ref{ct-uniform-generators}, with
base field $K_\psi$, and put
\[
 \epsilon_0=\frac{9\epsilon}{10},\quad
 R_*=R_2+\sum_{p\in P}R_p,\quad
 \eta=\frac{\epsilon_0}{2C_K},\quad
 J_0=2^{R_*+1}(401/\eta)^{R_*}.
\]
Fix any $J\ge J_0$ and any integer $h_0\ge h_*$.
Set $\widetilde j_p=\lfloor\log_pR_p\rfloor+3$.
There are arbitrarily large $Z$ for which one can choose a single
$U>1$ and a single pair $q,d$ satisfying
\eqref{sa-prime-ranges}, the congruences of
Theorem~\ref{qs-selection}, and $(qd,2M)=1$, such that the following
hold simultaneously for every $x\in[Z,Z^{3/2}]$.

There is an integer $h=h(x)\ge h_0$ with
\begin{equation}\label{sa-rank-window}
 x^{1/J}\le U\le x,\qquad
 x^{2+8\epsilon}\le c2^h\le2x^{2+8\epsilon}.
\end{equation}
All twists specified in Lemma~\ref{sa-character-families} with a
nonzero quadratic factor obey \eqref{qs-bounds} at $y=x$.
The conductor bounds \eqref{sa-conductor} hold, and for every
integer $1\le t\le x^{400}$,
\begin{equation}\label{sa-partial-cohomology}
 \sum_{1\le e<t}\lambda_2(E_h,e)
 +\sum_{p\in P}\sum_{1\le e<t}
              \lambda_p((K_p)_{\widetilde j_p},e)
       \le \exp(x^{4\epsilon/5}).
\end{equation}
All lengths in this formula are finite. The primes $q,d$ exceed
$2x^{401}+1$ and any further fixed bound specified in advance.
\end{proposition}

\begin{proof}
All constants in the statement are fixed before $X$ is chosen.
Apply Theorem~\ref{thm:unconditional-qs} at $U=X^{1/J}$ with its interval
exponent equal to $J$. The character bounds then hold throughout
$[U,U^J]=[U,X]$. By \eqref{sa-size-comparison},
\[
 \log(qd)\le X^{\epsilon_0/(4J)}
              \le X^{\epsilon_0/(4J_0)}.
\]
Proposition~\ref{ct-conditional-interval}, applied with
$\epsilon_0$, $C_t=401$, and $A=6$, gives $W$ with
$X^{1/J_0}\le W$, $W^2\le X$, and its cohomological estimates
on the whole interval $[W,W^2]$. In particular $W\ge U$.

Let $k=c^{1/\alpha}\ge1$, and for $s\ge1$ put
$e(s)=\lceil\alpha\log_2s\rceil$. Then
$s^\alpha\le2^{e(s)}\le2s^\alpha$.
Increase $X$ until $U\ge k^{400}$, $U\ge k^3$, and
$e(U)\ge h_0$. Set $Z=kW$. For $Z\le x\le Z^{3/2}$,
\[
 x\le(kW)^{3/2}\le W^2\le X,\qquad
 W\le x'=x/k\le W^2.
\]
Take $h=e(x')$. Then $h\ge h_0$ and
$x^\alpha=c(x')^\alpha\le c2^h\le2x^\alpha$.
Moreover $x^{1/J}\le X^{1/J}=U\le W\le x$.
Since $\alpha<3$ and $\log2>1/2$, one has
$h\le1+6\log x'\le6(1+\log x')$, as required by the
cohomological proposition. Finally
\[
                  x^{400}=k^{400}(x')^{400}\le(x')^{401}.
\]
Thus its estimate applies to every weight under consideration.

The layer used there is
$j_p=1+\lceil\log_pR_p\rceil$ (with $j_p=1$ for $R_p=1$).
One has $0\le\widetilde j_p-j_p\le2$.
Lemma~\ref{ct-growth} increases each odd-prime length by at
most $2R_p$ when passing to $\widetilde j_p$.
Divide the weighted estimate in
\eqref{ct-summed-bound} by $\log2$, the smallest of its prime
weights. Its left side dominates the corresponding unweighted
sum. After the layer changes the result is bounded by
\[
 \frac{\exp((x')^{18\epsilon/25})}{\log2}
       +2x^{400}\sum_{p\in P}R_p.
\]
Because $18/25<4/5$, this is at most
$\exp(x^{4\epsilon/5})$ for all sufficiently large $x$.
This proves \eqref{sa-partial-cohomology}, including finiteness.
The character and conductor conclusions follow from $[Z,Z^{3/2}]
\subset[U,X]$. Taking $X$ arbitrarily large makes $W$, hence $Z$,
arbitrarily large.
\end{proof}

The auxiliary exponent $\epsilon_0$ supplies a strict margin for the
fixed-layer change; it does not change the exponent in the prime
selection or in the rank window. The finite fields and generator bounds
are fixed before $J_0$. The potentially very large constant $J$ is
therefore an allowed fixed interval exponent in
Theorem~\ref{thm:unconditional-qs}.

\subsection{The full cohomological bound and the remaining input}

For a CM field $E$ and an odd prime $p$, let $\iota$ be the action
of complex conjugation on $\mathcal S_p(E,e)$. Its two parts relevant
here are the eigenspaces with eigenvalues $(-1)^e$ and $(-1)^{e+1}$;
write their lengths as $\lambda_p^+(E,e)$ and
$\lambda_p^-(E,e)$. Since $2$ is a unit in $\mathbb Z_p$, the
projections $(1\pm(-1)^e\iota)/2$ give a direct sum and hence
\begin{equation}\label{sa-length-decomposition}
       \lambda_p(E,e)=\lambda_p^+(E,e)+\lambda_p^-(E,e).
\end{equation}
This identity also holds for infinite lengths. Finiteness must be
established before ordinary finite sums of lengths are used.

The precise additional estimates needed to complete the
cohomological input are, for every finite set $S$ of odd primes,
\begin{align}\label{sa-missing-parts}
 \sum_{p\in S}(\log p)\sum_{1\le e<t}\lambda_p^+(E_h,e)
       &\le (N_ht)^C+C\exp(x^{4\epsilon/5}),\notag\\
 \sum_{p\in S}(\log p)\sum_{1\le e<t}\lambda_p^-(E_h,e)
       &\le K\exp(8x^{\epsilon/4}),
\end{align}
with fixed $C,K\ge0$ and with each individual length finite.
The source obtains the first from its arithmetic counting assertions
using the fixed odd-prime bounds, and the second from its separate
estimate for the other eigenspaces. The first line is not a statement
about only the finite set $P$: it controls every odd prime.
Neither line is proved by Proposition~\ref{sa-compatible-interval}.

\begin{corollary}\label{sa-full-bound}
If \eqref{sa-missing-parts} holds for the fields and scales in
Proposition~\ref{sa-compatible-interval}, then, for every finite
set $S$ of primes and $1\le t\le x^{400}$,
\[
 \sum_{p\in S}(\log p)\sum_{1\le e<t}\lambda_p(E_h,e)
       \le \sqrt{\exp(x^\epsilon)}
\]
for all sufficiently large $x$ in those intervals. Every individual
length occurring here is finite.
\end{corollary}
\begin{proof}
Use \eqref{sa-length-decomposition},
\eqref{sa-partial-cohomology}, and \eqref{sa-conductor}.
The left side is at most
\[
 (\log2+C)\exp(x^{4\epsilon/5})
       +\exp(2C x^{\epsilon/4})+K\exp(8x^{\epsilon/4}).
\]
Put $b=x^{4\epsilon/5}$, $s=x^{\epsilon/5}$,
$D=9+2C$, and $A=2+C+K$.
For $x\ge1$ the last display is at most $A\exp(Db)$.
If $s\ge2(D+\log A+1)$, then
$\log A+Db\le bs/2=x^\epsilon/2$, since $b\ge1$.
For example $x\ge[2(D+\log A+1)]^{5/\epsilon}$ suffices.
The stated finiteness follows from the hypotheses on the two
parts and the two-adic estimate.
\end{proof}

There is a useful separate check of the original rescaling. Suppose
one already has the full cohomological interval statement, for a fixed
function $e(s)$ satisfying
$s^\alpha\le2^{e(s)}\le2s^\alpha$, with weights below $s^{401}$
and bound $\sqrt{\exp(s^\epsilon)}$. Its quantifiers are:
a fixed $J$ first; then every sufficiently large $X$; then every
permissible prime pair at $U=X^{1/J}$; finally an interval
$[W,W^2]$ with $U\le W$, $W^2\le X$.
Select the prime pair on $[U,X]$ before invoking this statement.
The substitution $k=c^{1/\alpha}$, $Z=kW$, $s=x/k$ used above
then proves the full cohomological bound for weights below $x^{400}$
throughout $[Z,Z^{3/2}]$: the weight inclusion is
$x^{400}\le s^{401}$, and
$\sqrt{\exp(s^\epsilon)}\le\sqrt{\exp(x^\epsilon)}$.
This is exactly the role of the exponents $401$, $400$, and $3/2$.
No separate intersection of prime-selection and cohomological
intervals is taken.

One tempting change of scale is false. If a pair satisfies
\eqref{sa-prime-ranges} at $U$ and $x=U^J$ with fixed $J>1$,
then, for sufficiently large $U$,
\[
          \log q\le U^{\epsilon/5}
                <\tfrac12 x^{\epsilon/5}.
\]
Thus the pair need not satisfy the required lower window at $x$.
The construction retains $U$ as its selection scale and uses only
$x^{1/J}\le U\le x$, exactly as in
\eqref{sa-rank-window}. Likewise, the character family is fixed
through $\psi$; replacing it by all characters of the growing
$2$-power modulus would be a stronger, unproved selection assertion.


% Conventional reconstruction of the actual coefficient-lattice duality.
% All labels in this fragment have the prefix coef:.
